\section{Threats to Validity}
\label{sec:threats-to-validity}

\textbf{Dataset size and representativeness.}
The skills in our dataset may not cover all malicious skill families,
benign developer workflows, or future agent-skill ecosystems.
MalSkillBench and Datadog labels are inherited from their released
upstream partitions. In contrast, the VoltAgent samples are
source-benign skills randomly selected from a curated repository and
were not independently security-adjudicated by us. Residual benign-label
noise may therefore remain in this subset. Exact content hashing removes
exact duplicates but does not guarantee removal of semantically
equivalent or near-duplicate skills.
% The skills in our dataset may not cover all malicious skill families, benign developer workflows, or future agent-skill ecosystems. Labels are inherited from the source datasets and repositories and may therefore contain source-specific labeling noise. Exact content hashing removes exact duplicates but does not guarantee removal of semantically equivalent or near-duplicate skills.

\textbf{Sandbox coverage}
Dynamic analysis depends on whether sandbox execution triggers the relevant behavior. Some malicious skills may require specific user inputs, credentials, environment variables, external services, or delayed execution. A benign dynamic result does not prove the absence of malicious behavior. Safety-constrained blocking means that some evidence represents source-grounded dangerous intent rather than successful real-world harm.

\textbf{Run-to-run variability}
Each skill underwent a single sandbox execution.
The resulting dynamic evidence was saved and reused
in both the dynamic-evidence and Full Pipeline
configurations. Their comparison therefore uses
the same collected runtime evidence rather than
evidence from separate sandbox executions.
However, repeated sandbox executions and repeated
reasoning runs were not evaluated. The execution
agent may explore different actions across runs,
and the reasoning layer may interpret the same
evidence differently. The reported results therefore
describe the evaluated runs and do not establish
run-to-run stability.

\textbf{Human escalation and latency}
The architecture supports human escalation after three unresolved
reasoning rounds; however, no held-out test sample required human
adjudication. All predictions reported in
Tables~\ref{tab:rag_evaluation_metrics} and~\ref{tab:full_vs_baselines}
were produced automatically by the pipeline without human modification
of the final labels. This observation applies to the held-out
evaluation and does not imply that future ambiguous samples will never
require escalation. Latency values in
Table~\ref{tab:operational_latency} cover only the agentic stage and
should not be interpreted as end-to-end pipeline latency.
% The design supports human escalation after three unresolved reasoning rounds. The interpretation of Tables~\ref{tab:rag_evaluation_metrics} and~\ref{tab:full_vs_baselines} therefore depends on the escalation count and whether human judgments entered the reported predictions, as stated in Section~\ref{sec:operational-results}. Latency values in Table~\ref{tab:operational_latency} cover only the agentic stage and should not be interpreted as end-to-end pipeline latency.

\textbf{Adversarial resilience}
We do not evaluate skills adaptively engineered to evade our static rules, sandbox constraints, provenance logic, or semantic reasoning. The current results therefore measure performance on the benchmark corpus rather than worst-case robustness against an adaptive attacker.

\textbf{Scanner verdict interpretation.}
For Cisco Skill Scanner with LLM, we use its native \texttt{is\_safe} verdict rather than imposing a severity-based threshold. The native verdict does not always align with finding severity: both \texttt{wallet-tracker} and \texttt{emotion-bird} are malicious ground-truth samples, yet the former is marked unsafe with \textsc{Critical} severity, whereas the latter is marked safe despite multiple \textsc{Medium}-severity findings. Therefore, our reported metrics reflect the scanner's native binary decision; alternative severity-based rules could yield different classifications.